% Matrice impact × vraisemblance (grille pédagogique 1 à 4 de la notice)
\begin{tikzpicture}[x=9mm,y=7mm,font=\scriptsize]
  \foreach \v in {1,...,4}{\foreach \i in {1,...,4}{
    \pgfmathtruncatemacro\s{\v*\i}
    \ifnum\s>8 \def\c{alerte!22}\else\ifnum\s>4 \def\c{ambre!22}\else\def\c{vert!15}\fi\fi
    \fill[\c] (\v-1,\i-1) rectangle (\v,\i);
    \draw[white, line width=1pt] (\v-1,\i-1) rectangle (\v,\i);
    \node[gris, anchor=north east, font=\tiny] at (\v,\i) {\s};}}
  \newcommand\RisqueCase[3]{\node[font=\scriptsize\bfseries, align=center, text width=8mm] at (#1-0.5,#2-0.5) {#3};}
  \RisqueCase{3}{4}{R01, R02}
\RisqueCase{2}{4}{R03}
\RisqueCase{3}{3}{R04, R06}
\RisqueCase{2}{2}{R05}

  \foreach \v in {1,...,4}{\node at (\v-0.5,-0.3) {\v};}
  \foreach \i in {1,...,4}{\node at (-0.3,\i-0.5) {\i};}
  \node at (2,-0.85) {Vraisemblance};
  \node[rotate=90] at (-0.85,2) {Impact};
\end{tikzpicture}
